\chapter[1998 --- Furchgott et al.]{1998: Robert F. Furchgott, Louis J. Ignarro, Ferid Murad}
\label{ch:1998_Furchgott}

\section*{Main Contribution}

\textit{Discovered that nitric oxide acts as a signaling molecule in the cardiovascular system, explaining how blood vessels relax and leading to new treatments for heart disease.}\cite{nobel-1998,lecture-1998}

\section{Official Citation}

for their discoveries concerning nitric oxide as a signalling molecule in the cardiovascular system

\section{Background and Historical Context}

For most of the twentieth century, nitric oxide (NO) was regarded as a toxic environmental pollutant---a component of cigarette smoke and automobile exhaust, and a contributor to acid rain and ozone depletion. The idea that such a simple, reactive molecule could serve as a biological signaling agent seemed absurd. Yet that is precisely what Robert F. Furchgott, Louis J. Ignarro, and Ferid Murad independently discovered in the 1980s, a revelation that earned them the 1998 Nobel Prize in Physiology or Medicine and fundamentally transformed our understanding of cardiovascular physiology.

Robert Frederick Furchgott was born on June 4, 1916, in Charleston, South Carolina. He studied pharmacology at Cornell University, earning his Ph.D.\ in 1940, and spent his entire academic career at the State University of New York (SUNY) Downstate Medical Center in Brooklyn, where he became professor of pharmacology. Furchgott's research career spanned six decades and was characterized by meticulous experimental technique and a keen eye for anomalous observations.

Louis Joseph Ignarro was born on May 31, 1941, in Mobile, Alabama. He studied pharmacy at the University of Minnesota and earned his Ph.D.\ in pharmacology from the University of Minnesota in 1966. After postdoctoral work at the University of Minnesota and at the University of Virginia, Ignarro joined the faculty at the University of California, Los Angeles (UCLA), where he conducted his Nobel Prize--winning research on the identification of endothelium-derived relaxing factor (EDRF) as nitric oxide.

Ferid Murad was born on September 14, 1936, in Whiting, Indiana, to Albanian immigrant parents. He studied medicine at Indiana University and earned his M.D.\ and Ph.D.\ in pharmacology at the University of Minnesota in 1965. Murad joined the faculty at the University of Virginia and later moved to Stanford University, where he conducted his pioneering studies on the role of nitric oxide in cyclic GMP signaling. He later became chairman of the Department of Integrative Biology and Pharmacology at the University of Texas Health Science Center in Houston.

The intellectual context of their work was the discovery of cyclic GMP (cGMP) as a second messenger in the 1960s and 1970s. cGMP, like its better-known cousin cAMP, is a nucleotide that mediates the effects of hormones and other signaling molecules inside cells. By the late 1970s, it was known that cGMP levels were elevated in blood vessels treated with certain vasodilators, including acetylcholine, but the mechanism by which these agents stimulated cGMP production was unknown.

\section{The Discovery/Contribution}

\subsection{Furchgott and the Endothelium-Dependent Relaxation}

Furchgott's contribution began with an observation that initially seemed like an experimental artifact. In the early 1980s, he was studying the effects of acetylcholine on isolated blood vessel rings---a standard pharmacological preparation in which segments of artery are mounted in an organ bath and their contraction or relaxation is measured. Furchgott observed that acetylcholine sometimes caused relaxation of precontracted vessel rings and sometimes did not, depending on the preparation method.

Through careful experimentation, Furchgott discovered that the key variable was the endothelium---the single layer of cells lining the interior of the blood vessel. When the endothelium was intact, acetylcholine caused relaxation; when the endothelium was damaged or removed (by rubbing the interior of the vessel), acetylcholine had no relaxing effect and actually caused contraction. Furchgott concluded that the endothelium produced a short-lived, diffusible substance that he named ``endothelium-derived relaxing factor'' (EDRF), which caused the smooth muscle cells in the vessel wall to relax.

Furchgott published this finding in a landmark paper in 1980 and followed it with a series of papers characterizing the properties of EDRF. He showed that EDRF was highly unstable, with a half-life of only a few seconds, and that its action was mediated by the stimulation of guanylyl cyclase---the enzyme that converts GTP to cGMP---in smooth muscle cells.

\subsection{Ignarro and the Identification of EDRF as Nitric Oxide}

Independently, Ignarro was studying the same phenomenon from a different angle. In the early 1980s, he was investigating the role of cGMP in vasodilation and had shown that certain vasodilators stimulated guanylyl cyclase in smooth muscle cells. When he learned of Furchgott's discovery of EDRF, Ignarro conducted experiments comparing the chemical properties of EDRF with those of nitric oxide. He found that EDRF and nitric oxide had identical biological activities: both stimulated guanylyl cyclase, both caused vasodilation, and both were inactivated by hemoglobin and superoxide anion.

In 1986, Ignarro published a paper in \textit{Biochemical and Biophysical Research Communications} in which he proposed that EDRF was, in fact, nitric oxide. This identification was initially controversial, but it was soon confirmed by multiple laboratories, including Furchgott's own.

\subsection{Murad and the Role of NO in cGMP Signaling}

Murad's contribution preceded and complemented those of Furchgott and Ignarro. In the late 1970s, while at the University of Virginia, Murad had observed that nitric oxide and other nitrogen oxides stimulated guanylyl cyclase in tissue homogenates. He published these findings in 1977 and proposed that nitrogen oxides might play a role in the regulation of cGMP levels in cells. However, Murad did not initially connect his findings to the physiological phenomenon of endothelium-dependent vasodilation.

After the identification of EDRF as nitric oxide, Murad's earlier work was recognized as having anticipated the key insight. The convergence of Furchgott's pharmacological observations, Ignarro's chemical identification, and Murad's biochemical characterization of guanylyl cyclase activation provided a complete and compelling picture of nitric oxide as a biological signaling molecule.

\subsection{The NO-cGMP Pathway}

The combined work of the three laureates revealed a signaling pathway of remarkable elegance. In blood vessels, the neurotransmitter acetylcholine (or other endothelium-dependent stimuli) activates endothelial cells, which produce nitric oxide from the amino acid L-arginine by the action of the enzyme endothelial nitric oxide synthase (eNOS). The nitric oxide diffuses across the short distance from the endothelium to the underlying smooth muscle cells, where it activates soluble guanylyl cyclase, stimulating the production of cGMP. cGMP activates protein kinase G, which phosphorylates target proteins that cause smooth muscle relaxation and vasodilation.

\section{Impact on Modern Medicine}

The discovery of nitric oxide as a signaling molecule has had an enormous impact on cardiovascular medicine and beyond.

\subsection{Cardiovascular Disease}

Endothelial dysfunction---the impaired production or bioavailability of nitric oxide---is now recognized as a key early event in the development of atherosclerosis, hypertension, and heart failure. Risk factors for cardiovascular disease, including smoking, diabetes, high cholesterol, and aging, all reduce the bioavailability of nitric oxide. Strategies to restore or enhance nitric oxide signaling, including the use of statins, ACE inhibitors, and lifestyle modifications, are central to the prevention and treatment of cardiovascular disease.

\subsection{Phosphodiesterase Inhibitors and Erectile Dysfunction}

The NO-cGMP pathway is the mechanism of action of the phosphodiesterase type 5 (PDE5) inhibitors---sildenafil (Viagra), tadalafil (Cialis), and vardenafil (Levitra)---which are used to treat erectile dysfunction. These drugs inhibit the enzyme that degrades cGMP, thereby prolonging and enhancing the vasodilatory effects of nitric oxide in the corpus cavernosum. The discovery that NO is the physiological mediator of penile erection was a direct consequence of the work of Furchgott, Ignarro, and Murad.

\subsection{Pulmonary Hypertension}

Nitric oxide is also used therapeutically in the treatment of pulmonary hypertension. Inhaled nitric oxide selectively dilates the pulmonary vasculature, reducing pulmonary arterial pressure and improving oxygenation. This application, approved by the FDA in 1999, is based directly on the signaling pathway elucidated by the three laureates.

\subsection{Platelet Aggregation and Thrombosis}

Nitric oxide inhibits platelet aggregation and adhesion, reducing the risk of thrombosis. This antithrombotic property of NO is mediated by the same cGMP-dependent mechanism that causes vasodilation and is an important component of the endothelium's protective function.

\subsection{Neurotransmission}

Nitric oxide also functions as a neurotransmitter in the central and peripheral nervous systems, playing roles in learning, memory, and gastrointestinal motility. The recognition of NO as a unconventional neurotransmitter---one that is not stored in vesicles but is synthesized on demand and diffuses freely across cell membranes---has expanded our understanding of neural signaling.

\section{Legacy}

Robert Furchgott continued his research at SUNY Downstate Medical Center until his retirement. He received numerous honors, including the Albert Lasker Basic Medical Research Award in 1996. Furchgott passed away on May 19, 2009, at the age of 92.

Louis Ignarro became professor emeritus of molecular and medical pharmacology at UCLA and has been a prominent advocate for the health benefits of nitric oxide, including the nutritional supplement Neo40. He received numerous honors and continues to lecture widely on the role of NO in health and disease.

Ferid Murad became director of the Institute for Molecular Medicine at the University of Texas Health Science Center in Houston and continued to study the role of cyclic nucleotides in cell signaling. He received numerous honors and passed away on September 4, 2023, at the age of 86.

Together, Furchgott, Ignarro, and Murad revealed that one of the simplest molecules in chemistry---nitric oxide, a gas with a single unpaired electron---serves as one of the most versatile signaling molecules in biology. Their discovery has transformed our understanding of cardiovascular physiology and has led to the development of drugs that have improved the lives of millions of patients worldwide.


\section{Modern Developments}

Furchgott, Ignarro, and Murad's nitric oxide work has led to modern cardiovascular medicine. Understanding of NO signaling has led to treatments for pulmonary hypertension (sildenafil, originally for erectile dysfunction). Understanding of endothelial function has led to therapies for heart failure (nesiritide). Understanding of NO in immunity has revealed its role in macrophage antimicrobial defense. Understanding of NO in the nervous system has informed development of treatments for neurological disorders.


\section{Bibliography}

\begin{thebibliography}{9}
\bibitem{nobel-1998}
Nobel Prize Outreach, ``The Nobel Prize in Physiology or Medicine 1998,''\emph{NobelPrize.org}.\\
\url{https://www.nobelprize.org/prizes/medicine/1998/summary/}

\bibitem{lecture-1998}
Robert F. Furchgott, ``Nobel Lecture,''\emph{NobelPrize.org}. Winner-authored primary source; downloadable from the official Nobel Prize archive.\\
\url{https://www.nobelprize.org/prizes/medicine/1998/furchgott/lecture/}
\end{thebibliography}
